\documentclass[10pt, a4paper]{article}
\usepackage{preamble}

\title{Data Science and Statistical Modelling II \\
    \large Cheat Sheet}
\author{Luke Phillips}
\date{March 2026}

\begin{document}

\maketitle

\newpage

\section{Computational Statistics}

\textbf{Monte Carlo testing}:
Given data $(x_1, \dotsc, x_n)$,
an observed test statistic $t = h(x_1, \dotsc, x_n)$,
a data generating distribution $F(x \mid \theta)$ and a hypothesis,
$H_0 : \theta = \theta_0$,
$H_1 : \theta > \theta_0$.
We do the following steps to perform a Monte Carlo hypothesis test:
For $j \in \{1, \dotsc, N\}$
\begin{enumerate}[label = (\arabic*)]
    \item Simulate $n$ observations $(z_1, \dotsc, z_n)$ where $Z_i \sim F(\cdot\mid\theta_0)$.

    \item Calculate $t_j = h(z_1, \dotsc, z_n)$.
\end{enumerate}
Then estimate the $p$-value of the test
\[
\P(T \geq t \mid H_0\text{ true}) \approx \frac{1}{N}\sum_{j = 1}^{N}\mathbbm{1}\{t_j \geq t\}
\]


\textbf{Bootstrap}:
Given a data set consisting of independent samples $\vec{x} = (x_1, \dotsc, x_n)$ and statistic $S(\cdot)$.
To construct a bootstrap estimate of the standard error of the statistic $S(\cdot)$,
draw $B$ new samples of size $n$ with replacement from $\vec{x} = (x_1, \dotsc, x_n)$,
giving $\vec{x} ^ {*1}, \dotsc, \vec{x} ^ {*B}$ and compute
\[
\widehat{\Var}(S(\vec{x})) = \frac{1}{B - 1}\sum_{b = 1}^{B}(S(\vec{x} ^ {*b}) - \bar{S ^ {*}}) ^ 2
\]
where $\bar{S ^ {*}} = \frac{1}{B}\sum_{b = 1}^{B}S(\vec{x} ^ {*b})$.

\textbf{The ECDF}
The ecdf is defined as
\[
\hat{F}(x) \coloneqq \frac{1}{n}\sum_{i = 1}^{n}\mathbbm{1}\{x_i \leq x\}\qquad\forall x \in \R.
\]
For given sample data $(x_1, \dotsc, x_n)$.

For a finite population size $N$,
sample mean $\bar{X}$ based on sample size $n$ uniformly drawn without replacement
\[
\Var(\bar{X}) = \left(\frac{N - n}{N - 1}\right)\frac{\sigma ^ 2}{n}.
\]

\textbf{Population bootstrap}
Assume $\vec{x} = (x_1, x_2, \dotsc, x_n)$ is a sample of data drawn from a finite population size $N$,
$S(\cdot)$ the statistic we wish to use as an estimator.
Further assume $N / n = k$ is an integer.

To perform the \textbf{population bootstrap},
construct a pseudo-population of size $N$ by repeating the samples $\vec{x}$ $k$ times:
\[
\tilde{\vec{x}} = (\underbrace{x_1, \dotsc, x_n}_{\vec{x}}, \underbrace{x_1, \dotsc, x_n}_{\vec{x}}, \dotsc, \underbrace{x_1, \dotsc, x_n}_{\vec{x}}).
\]
Then,
replicate the finite population sampling scheme by drawing $B$ new samples of size $n$ without replacement from $\tilde{\vec{x}}$,
giving $\vec{x} ^ {*1}, \dotsc, \vec{x} ^ {*B}$ then proceed as for the standard bootstrap:
\[
\widehat{\Var}(S(\vec{x})) = \frac{1}{B - 1}\sum_{b = 1}^{B}(S(\vec{x} ^ {*b}) - \bar{S} ^ {*}) ^ 2
\]
where $\bar{S} ^ {*} = \frac{1}{B}\sum_{b = 1}^{B}S(\vec{x} ^ {*b})$.
If $n \nmid N$,
that is $N = kn + m$ for $0 < m < n$,
then before each bootstrap sample form pseudo-population $\tilde{\vec{x}}$ by taking $k$ copies of $\vec{x}$ and add an additional subsample of size $m$,
without replacement,
from $\vec{x}$.


\end{document}